\documentclass[12pt]{article}
\usepackage[brazil]{babel}
\usepackage[utf8]{inputenc}
\usepackage[T1]{fontenc}
\usepackage{amsmath, amssymb, amsthm}
\usepackage{geometry}
\usepackage{graphicx}
\usepackage{xcolor}
\usepackage{enumitem}
\usepackage{listings}

\definecolor{maincolor}{RGB}{0,100,50}
\usepackage{fancyhdr}
\definecolor{forestgreen}{RGB}{0,100,50}
\usepackage[most]{tcolorbox}

\lstset{
  basicstyle=\ttfamily\small,
  keywordstyle=\color{blue},
  stringstyle=\color{red},
  commentstyle=\color{green!50!black},
  breaklines=true,
  frame=single,
  numbers=left,
  numberstyle=\tiny\color{gray},
}

\newtcolorbox{block}[1]{%
  colback=green!5,
  colframe=maincolor,
  coltitle=white, 
  title=#1,
  fonttitle=\bfseries,
  boxrule=1.5pt,
  arc=3mm,
  left=6pt,
  right=6pt,
  top=6pt,
  bottom=6pt
}

\geometry{margin=2.5cm}
\pagestyle{fancy}
\fancyhf{}

\renewcommand{\headrulewidth}{2pt}
\renewcommand{\footrulewidth}{2pt}

\renewcommand{\headrule}{\hbox to\headwidth{%
  \color{forestgreen}\leaders\hrule height \headrulewidth\hfill}}

\renewcommand{\footrule}{\hbox to\headwidth{%
  \color{forestgreen}\leaders\hrule height \footrulewidth\hfill}}

\setlength{\headheight}{32pt}

\fancyhead[L]{%
  \includegraphics[height=1.2cm]{Universidade_de_Rio_Verde.jpg}
}

\fancyhead[R]{%
  \raggedleft
  \textcolor{forestgreen}{%
    \small
    \textbf{Atividade de Programação.}\\
    Me.~Sergio Souza Novak\\
    \textit{Lógica para Engenharia de Software}
  }
}

\fancyfoot[C]{%
  \textcolor{forestgreen}{\thepage}
}

\begin{document}

\vspace*{0.2cm}
\section*{Exercícios de Programação — Relações Binárias}

\begin{block}{Instruções Gerais}
Para cada um dos exercícios abaixo, escolha uma das seguintes linguagens de programação: \textbf{JavaScript, Java, C/C++ ou Python}. 

\textbf{Regras de Entrada:}
\begin{itemize}
    \item O programa deve ler dois conjuntos, $A$ e $B$, a partir do teclado.
    \item Os elementos dos conjuntos devem ser \textbf{números naturais} ($0, 1, 2, \dots$).
    \item A leitura de cada conjunto deve ser interrompida quando o usuário digitar um \textbf{número menor que zero} (ex: -1).
\end{itemize}

\textbf{Formato de Entrega:}
\begin{itemize}
    \item A entrega deve ser um arquivo \textbf{.zip}.
    \item Dentro do arquivo .zip, deve haver \textbf{uma pasta para cada exercício} (Ex: Exercicio1, Exercicio2, etc.).
    \item Dentro de cada pasta, devem estar contidos \textbf{todos os códigos e arquivos necessários} para compilar e executar a resposta.
\end{itemize}
\end{block}

\vspace{0.4cm}
\begin{enumerate}

\item \textbf{Produto Cartesiano}\\
Implemente um programa que receba os conjuntos $A$ e $B$ e apresente na tela o produto cartesiano $A \times B$. O resultado deve ser exibido como um conjunto de pares ordenados, por exemplo: $\{(a_1, b_1), (a_1, b_2), \dots\}$.

\begin{block}{Exemplo}
\textbf{Entrada:} \\
Conjunto A: 1, 2, -1 (fim) \\
Conjunto B: 3, 4, -1 (fim) \\
\textbf{Saída:} \\
$A \times B = \{(1, 3), (1, 4), (2, 3), (2, 4)\}$
\end{block}

\vspace{0.3cm}
\item \textbf{União de Produtos Cartesianos}\\
Implemente um programa que calcule e mostre a união dos produtos cartesianos $A \times B$ e $B \times A$. Ou seja, o conjunto:
\[ (A \times B) \cup (B \times A) \]

\begin{block}{Exemplo}
\textbf{Entrada:} \\
Conjunto A: 1, -1 \\
Conjunto B: 2, -1 \\
\textbf{Saída:} \\
$(A \times B) \cup (B \times A) = \{(1, 2), (2, 1)\}$
\end{block}

\vspace{0.3cm}
\item \textbf{Propriedades de uma Relação}\\
Dada uma relação $R$ lida do teclado (onde os alunos devem definir como os pares $(x, y)$ da relação serão informados, utilizando os elementos dos conjuntos $A$ e $B$ previamente lidos), verifique as seguintes propriedades:
\begin{enumerate}[label=\alph*)]
    \item \textbf{Reflexividade:} Verifique se a relação é reflexiva em $A$.
    \item \textbf{Simetria:} Verifique se a relação é simétrica.
    \item \textbf{Antissimetria:} Verifique se a relação é antissimétrica.
    \item \textbf{Assimétrica:} Verifique se a relação é assimétrica.
\end{enumerate}

\begin{block}{Exemplo}
\textbf{Entrada:} \\
Conjunto A: 1, 2, -1 \\
Relação R: (1, 1), (2, 2), (1, 2), -1 (fim da relação) \\
\textbf{Saída:} \\
Reflexiva: Sim \\
Simétrica: Não (falta (2, 1)) \\
Antissimétrica: Sim \\
Assimétrica: Não (contém (1, 1))
\end{block}

\vspace{0.3cm}
\item \textbf{Matriz de Relação}\\
Dada uma relação qualquer $R$ sobre os conjuntos lidos, implemente um código que gere e exiba a \textbf{Matriz de Relação} correspondente. A matriz deve ter as dimensões adequadas e exibir $1$ se existe a relação entre os elementos e $0$ caso contrário.

\begin{block}{Exemplo}
\textbf{Entrada:} \\
Conjunto A: 1, 2, -1 \\
Conjunto B: 1, 2, -1 \\
Relação R: (1, 2), (2, 1), -1 \\
\textbf{Saída:} \\
$\begin{pmatrix} 0 & 1 \\ 1 & 0 \end{pmatrix}$
\end{block}

\end{enumerate}

\end{document}
